\section{Background}
\label{sec:background}

\subsection{ArkTS Programs and Platform APIs}

\openharmony applications are organized around abilities, pages, ArkUI components, lifecycle methods, and event callbacks~\cite{openharmony,arktsdoc}. ArkTS retains TypeScript-like classes, functions, fields, closures, and Promises while adding a declarative UI model. A component's \texttt{build} method describes UI structure; methods such as \texttt{aboutToAppear} and handlers such as \texttt{onClick} are invoked by the framework. Privacy-relevant behavior can consequently span an initializer, a framework callback, component state, and a later event even when no explicit source-level call connects these methods.

ArkAnalyzer lowers ArkTS into ArkIR and represents anonymous functions, static and instance initialization, closures, exception flow, and access paths in generated or transformed methods~\cite{arkanalyzer}. These methods remain semantically relevant: they contain callback assignments, factory calls, captured values, and component initialization. Our analysis uses ArkAnalyzer's Scene, typed IR, call graph, and pointer evidence rather than interpreting ArkTS text in isolation.

Platform APIs add a second form of indirection. OpenHarmony has evolved from fine-grained \texttt{@ohos.*} modules toward consolidated \texttt{@kit.*} exports, and applications may mix both forms. An imported namespace can be renamed, assigned to an IR temporary, or used to create a typed manager on which the privacy operation is invoked. Some configured operations are executable property reads rather than calls. The analysis must therefore preserve API ownership across source syntax, SDK declarations, and ArkIR.

\subsection{Taint Analysis and Source Binding}

An interprocedural taint analysis starts from a set of source facts and propagates them through program statements until they reach configured sinks. IFDS represents this computation over an exploded interprocedural supergraph and supports precise, distributive flow functions~\cite{reps95,sagiv96}. HapFlow instantiates this approach for ArkTS and models closures, lifecycle entries, and source--sink queries on top of ArkAnalyzer~\cite{hapflow}.

Before propagation begins, each source specification must answer two questions. The \emph{identity} question determines whether a program operation denotes the configured platform API. The \emph{carrier} question determines which program value receives the modeled data. For a direct-return API, the carrier is normally the assigned result. For a callback API, it is a designated callback parameter. For a Promise API, it is the fulfilled value observed by a success continuation and, after \texttt{await}, the assigned result. An operation may be privacy-relevant without producing modeled data; such an occurrence remains detector evidence but does not create an IFDS seed.

We use \emph{source binding} for the accepted association between an API identity and its carrier. A source binding also retains a witness explaining why the owner and carrier were accepted. This evidence is important in ArkTS because package presence, member spelling, function type, and control-flow reachability are individually insufficient. Section~\ref{sec:design} defines the binding and shows how its witness is preserved during propagation.

\subsection{Scope of the Analysis}

\mytool answers two audit questions: which configured privacy-sensitive platform operations occur, and which modeled values may reach configured sinks. It does not equate an API occurrence or a may-flow with a policy violation. The detector and taint query consequently remain separate. Detector rules cover configured calls and executable properties; query rules additionally require a concrete return, callback, argument, or framework-input carrier. Detector-local sink observations provide nearby review context, whereas configured-query paths are produced by IFDS or explicitly labeled continuation recovery. Keeping these evidence sets distinct avoids interpreting every API occurrence as a leak or every nearby sink as a data-flow endpoint.
